\section{Динамические игры с полной информацией}

В данном разделе рассматривается класс игр, в которых участники принимают решения не одновременно, а последовательно. Переход от статических моделей к динамическим требует введения новых концепций, описывающих развернутую во времени структуру взаимодействия.

Основные характеристики динамической игры:
\begin{itemize}
    \item \textbf{Очередность ходов:} игроки совершают действия в строго определенной последовательности.
    \item \textbf{Неотвратимость выбора:} однажды сделанный ход не может быть отменен.
    \item \textbf{Наблюдаемость:} в играх с совершенной информацией игроки видят действия, совершенные предшественниками.
    \item \textbf{Терминальные платежи:} выигрыши распределяются только после завершения игры (в терминальных вершинах дерева).
\end{itemize}

\subsection{Концепция стратегии в динамической игре}

Рассмотрим классическую игру «Лобовая атака», в которой два водителя несутся навстречу друг другу. Каждый может либо свернуть (\textit{T}), либо продолжить движение (\textit{N}). 

\begin{example}[Лобовая атака: от статики к динамике]
В статической постановке (когда решения принимаются одновременно) матрица выигрышей выглядит следующим образом:

\begin{table}[h]
\centering
\renewcommand{\arraystretch}{1.3}
\begin{NiceTabular}{cc|cc}
\toprule
& & \multicolumn{2}{c}{\textbf{Михаэль}} \\
& & \textit{T (Свернуть)} & \textit{N (Не сворачивать)} \\
\midrule
\multirow{2}{*}{\textbf{Иван}} 
& \textit{T (Свернуть)} & $0, 0$ & $-5, 10$ \\
& \textit{N (Не сворачивать)} & $10, -5$ & $-10, -10$ \\
\bottomrule
\end{NiceTabular}
\end{table}

В статической игре существуют два равновесия Нэша: $(T, N)$ и $(N, T)$. Однако предположим, что один из водителей (Иван) обладает лучшей реакцией и успевает сделать ход первым, тогда как Михаэль принимает решение, уже наблюдая за действиями Ивана. 

Для решения такой динамической игры применяется метод \textbf{обратной индукции}. Анализ начинается с конца:
\begin{enumerate}
    \item Если Иван выбрал \textit{N}, Михаэль стоит перед выбором: \textit{T} (выигрыш $-5$) или \textit{N} (выигрыш $-10$). Рациональным выбором будет \textit{T}.
    \item Если Иван выбрал \textit{T}, Михаэль выбирает между \textit{T} (выигрыш $0$) и \textit{N} (выигрыш $10$). Рациональным выбором будет \textit{N}.
    \item Зная реакцию Михаэля, Иван на первом шаге сравнивает свои итоговые выигрыши: выбор \textit{N} принесет ему $10$, а выбор \textit{T} принесет $-5$. Следовательно, Иван выберет \textit{N}.
\end{enumerate}
\end{example}

Важнейшим аспектом динамических игр является строгое понимание термина «стратегия». Стратегия — это полный план действий, содержащий инструкции для \textit{каждой} возможной ситуации, в которой игроку может потребоваться сделать ход, даже если эта ситуация никогда не наступит в реальной партии.

В примере выше у Михаэля есть $2 \times 2 = 4$ возможные чистые стратегии:
\begin{itemize}
    \item \textbf{TT:} Всегда сворачивать.
    \item \textbf{TN:} Сворачивать, только если Иван не свернул.
    \item \textbf{NT:} Сворачивать, только если Иван свернул.
    \item \textbf{NN:} Никогда не сворачивать.
\end{itemize}

Нормальная форма данной динамической игры будет представлена матрицей $2 \times 4$:

\begin{table}[h]
\centering
\renewcommand{\arraystretch}{1.3}
\begin{NiceTabular}{cc|cccc}
\toprule
& & \multicolumn{4}{c}{\textbf{Михаэль}} \\
& & TT & TN & NT & NN \\
\midrule
\multirow{2}{*}{\textbf{Иван}} 
& T & $0, 0$ & $-5, 10$ & $0, 0$ & $-5, 10$ \\
& N & $10, -5$ & $10, -5$ & $-10, -10$ & $-10, -10$ \\
\bottomrule
\end{NiceTabular}
\end{table}

\subsection{Формализация игры в развернутой форме}

Игра в развернутой форме задается в виде ориентированного графа (дерева игры), где:
\begin{itemize}
    \item Каждой нетерминальной (внутренней) вершине соответствует игрок, делающий ход.
    \item Существует фиктивный игрок «Природа», совершающий ходы с заданной экзогенной вероятностью.
    \item Каждой терминальной (конечной) вершине соответствует вектор выигрышей всех игроков.
    \item Вершины объединяются в \textbf{информационные множества}. Игрок знает, в каком информационном множестве он находится, но не знает, в какой конкретно вершине внутри этого множества (что моделирует ненаблюдаемость чужих ходов или забывчивость).
\end{itemize}

\begin{definition}[Чистая стратегия в динамической игре]\label{def:pure_dyn}
В игре $\Gamma$ в развернутой форме множество чистых стратегий игрока $i$ определяется как декартово произведение:
\begin{equation}
    S_i = \prod_{h_i \in H_i} A(h_i)
\end{equation}
где $H_i$ — множество всех информационных множеств, в которых делает ход игрок $i$, а $A(h_i)$ — множество действий, доступных игроку $i$ в информационном множестве $h_i$. 

Общее число чистых стратегий игрока $i$ равно:
\begin{equation}
    N_i = \prod_{h_i \in H_i} |A(h_i)|
\end{equation}
\end{definition}

\subsection{Игры с совершенной информацией и Теорема Куна}

Игра называется игрой с \textit{совершенной информацией}, если каждое информационное множество состоит ровно из одной вершины (все игроки наблюдают всю предысторию игры).

\begin{theorem}[Кун, 1953]\label{thm:kuhn_perfect}
В любой конечной игре с совершенной информацией существует равновесие Нэша в чистых стратегиях.
\end{theorem}

\begin{proof}
Доказательство строится по индукции на основе алгоритма Цермело. Пусть $\Gamma$ — динамическая конечная игра. Рассмотрим семейство игр $\Gamma_1, \dots, \Gamma_T$, где $\Gamma_1 = \Gamma$. Игра $\Gamma_{t+1}$ получается из $\Gamma_t$ следующим образом:
\begin{enumerate}
    \item Выбирается предпоследняя вершина $x \in \Gamma_t$ (такая, что все следующие за ней вершины $Z = \{z_1, \dots, z_M\}$ являются терминальными). Пусть в $x$ ходит игрок $i$.
    \item В вершине $z^* \in Z$ выигрыш игрока $i$ максимален. Пусть $s(x)$ — действие, ведущее в $z^*$.
    \item Вершины $Z$ удаляются из дерева. Вершина $x$ объявляется терминальной с вектором выигрышей $u(x) = u(z^*)$.
\end{enumerate}
Процедура повторяется до тех пор, пока не останется только начальная вершина дерева. Поскольку на каждом шаге игрок максимизирует свой выигрыш и игра конечна, полученный набор действий формирует профиль чистых стратегий, от которого никому не выгодно отклоняться.
\end{proof}

\subsection{Равновесие, совершенное по подыграм}

Стандартное равновесие Нэша в динамических играх может порождать неправдоподобные исходы, основанные на пустых угрозах. Для исключения таких ситуаций вводится концепция подыгры.

\begin{definition}[Подыгра]
Пусть $\Gamma$ — игра в развернутой форме. Подыгра $\Gamma(x)$ — это часть дерева игры, которая начинается с некоторой вершины $x$ и удовлетворяет следующим свойствам:
\begin{itemize}
    \item Вершина $x$ является единственным элементом в своем информационном множестве.
    \item Информационные множества, содержащие вершины из $\Gamma(x)$, не содержат вершин, лежащих вне дерева подыгры $\Gamma(x)$ (подыгра не "разрезает" информационные множества).
\end{itemize}
\end{definition}

\begin{definition}[Равновесие, совершенное по подыграм -- SPE]
Профиль стратегий называется равновесием, совершенным по подыграм (Subgame Perfect Equilibrium), если он индуцирует равновесие Нэша в каждой подыгре исходной игры.
\end{definition}

\begin{example}[Вход на рынок]
Действующий монополист (Старожил) сталкивается с потенциальным конкурентом (Новичок). Новичок решает: \textit{Войти} или \textit{Нет}. Если он входит, Старожил решает: \textit{Подвинуться} (разделить рынок) или развязать \textit{Ценовую войну}.

\begin{table}[h]
\centering
\renewcommand{\arraystretch}{1.3}
\begin{NiceTabular}{cc|cc}
\toprule
& & \multicolumn{2}{c}{\textbf{Старожил}} \\
& & \textit{Подвинуться} & \textit{Война} \\
\midrule
\multirow{2}{*}{\textbf{Новичок}} 
& \textit{Войти} & $2, 1$ & $0, 0$ \\
& \textit{Нет} & $1, 2$ & $1, 2$ \\
\bottomrule
\end{NiceTabular}
\end{table}

В данной игре два равновесия Нэша: (\textit{Войти}, \textit{Подвинуться}) и (\textit{Нет}, \textit{Война}). Однако равновесие (\textit{Нет}, \textit{Война}) содержит неправдоподобную угрозу: если Новичок всё же войдет на рынок, Старожилу будет невыгодно начинать войну (выигрыш $0$ против $1$). Применив обратную индукцию, находим единственное SPE: (\textit{Войти}, \textit{Подвинуться}).
\end{example}

\subsection{Поведенческие стратегии}

Задание смешанной стратегии как распределения вероятностей на всем множестве чистых стратегий в динамических играх крайне громоздко. Удобнее использовать поведенческие стратегии.

Поведенческая стратегия формализует распределение вероятностей локально — в каждом информационном множестве:
\begin{equation}
    b_i(a_i | h_i) = \frac{\sum_{s_i \in R_i(h_i), s_i(h_i)=a_i} \sigma_i(s_i)}{\sum_{s_i \in R_i(h_i)} \sigma_i(s_i)}
\end{equation}
где $R_i(h_i)$ — подмножество чистых стратегий игрока $i$, допускающих прохождение через информационное множество $h_i$, а $\sigma_i(s_i)$ — вероятность сыграть чистую стратегию $s_i$.

\begin{theorem}[Кун, 1953]\label{thm:kuhn_behavioral}
Для всех игр с совершенной памятью смешанные и поведенческие стратегии эквивалентны.
\end{theorem}

\subsection{Игры с несовершенной информацией: Пример «Верю - не верю»}

Рассмотрим игру с участием Природы, порождающую скрытую информацию. 
У Игрока 1 на руках может быть карта: либо $6$, либо Туз (с равными вероятностями). Игрок 2 карту не видит. 
Игрок 1 делает заявление о своей карте. Если он заявляет «Туз», Игрок 2 может выбрать: \textit{Верю} (B) или \textit{Не верю} (N).

\textbf{Стратегии игроков:}
\begin{itemize}
    \item Игрок 1 задает вероятности $(p_1, p_2)$ — вероятность заявить «Туз», если на руках Туз ($p_1$), и вероятность заявить «Туз», если на руках $6$ ($p_2$).
    \item Игрок 2 задает вероятность $q$ — сказать «Не верю», если Игрок 1 заявил «Туз».
\end{itemize}

Заметим, что стратегия $p_1 < 1$ строго доминируется. Игроку 1 всегда выгодно заявлять «Туз», если у него на руках Туз. Таким образом, $p_1 = 1$. Ожидаемый выигрыш Игрока 1 от блефа зависит от $p_2$ и $q$:
\begin{equation}
    u_1(p_2, q) = -\frac{1-p_2}{2} - 2q \frac{p_2}{2} + (1-q)\frac{p_2}{2} + 2\frac{1}{2}q + \frac{1}{2}(1-q)
\end{equation}
После алгебраических упрощений функция полезности принимает вид:
\begin{equation}
    u_1(p_2, q) = p_2 - \frac{3}{2}p_2 q + \frac{q}{2}
\end{equation}

Для нахождения равновесия строятся функции наилучшего ответа (реакции):
\begin{equation}
    \check{p}_2(q) = 
    \begin{cases} 
    1, & \text{если } q < \frac{2}{3} \\ 
    [0, 1], & \text{если } q = \frac{2}{3} \\ 
    0, & \text{если } q > \frac{2}{3} 
    \end{cases}
\end{equation}

\begin{equation}
    \check{q}(p_2) = 
    \begin{cases} 
    1, & \text{если } p_2 > \frac{1}{3} \\ 
    [0, 1], & \text{если } p_2 = \frac{1}{3} \\ 
    0, & \text{если } p_2 < \frac{1}{3} 
    \end{cases}
\end{equation}

В точке пересечения данных функций находится равновесие Нэша в смешанных (поведенческих) стратегиях:
\begin{equation}
    p_2^* = \frac{1}{3}, \quad q^* = \frac{2}{3}
\end{equation}
В равновесии Игрок 1 блефует с вероятностью $1/3$, а Игрок 2 проверяет его с вероятностью $2/3$.
